\chapter{研究背景与阶段目标}\label{chap:middle-background}
\section{海缆巡检任务及研究意义}

海底电缆承担海上风电场内部集电与海陆之间的电能传输，其运行状态同时受到电气负荷、海床变化和外部机械作用的影响。路由偏移、冲刷裸露及局部悬跨等几何变化可能先于明显电气故障出现。因此，运维所需的信息不仅是电缆是否通电，还包括目标位于何处、观测是否可信，以及哪些区段需要进一步复核。针对海底电缆的损伤与寿命分析也强调，机械载荷与环境条件必须共同进入评估过程\cite{dinmohammadi2019predicting}。

自主水下机器人（autonomous underwater vehicle，AUV）能够在通信受限时执行预设任务，为沿路由持续采集数据提供载体。这也将目标识别、运动控制和异常处置的责任前移到艇载系统：航迹需要与目标关联，观测质量需要支持当前估计，命令需要符合近底运动约束。因此，海缆巡检要求联合研究感知、估计、决策与执行。

本课题面向存量 AUV 平台的改造，研究范围以控制软件、多源信息接入、磁感知和安全自主巡检为主。原艇结构、既有执行机构及 VxWorks 控制系统构成工程基础；声呐载荷的专项研制与加工由协作工作承担。本文需要解决的是如何将已有平台与新增感知计算能力衔接，并用分层实验说明这些衔接在何种条件下有效。研究并不以重新设计完整航行器作为前提，而是围绕可观测信息如何变成可执行、可追溯的巡检决策展开。

本文将 AUV 定位于沿缆巡查和异常筛查，精细接触检查与维修仍交由相应作业平台处理。系统应形成可追溯的位置、质量和状态记录，在信息不足时有明确的退出和转交路径。

\section{相关研究与本课题问题}

声学与磁学为目标提供不同类型的观测。侧扫声呐关注海床表面纹理、目标形态及声影，可为裸露或具有可识别表面特征的管缆提供方向和横向线索；浅地层剖面仪利用沉积层内部反射结构探查浅埋目标。二者的频段、成像几何与适用条件不同，不能仅以“声呐探测”合并其能力。已有 AUV 侧扫跟踪研究说明，声学目标提取能够进入自主制导流程\cite{bagnitsky2011sidescan,feng2023automatic}；本课题复用这种声学线索的接口思想，重点研究其与磁观测、导航状态及安全决策的关系。

磁观测与载流导体的电流和几何关系直接相关，能够补充依赖表面可见性的声学信息。已有磁导引及标量磁测跟踪研究分别探讨了利用磁场线索完成路由跟随以及降低姿态影响的途径\cite{xiang2016subsea,li2024attitude}。实际海缆中的多芯结构、相位抵消、铠装影响与邻近磁源，使磁场幅度同时依赖多个未知量。即使单条测线具有清晰峰值，也不能直接认定其对应指定电缆，更不能把某一理想场形下的宽度系数无条件换算为埋深。

本研究选择半高全宽（full width at half maximum，FWHM）作为主动横切剖面的可解释特征。其价值在于将空间幅度变化转成可检查的峰值、半高阈值和两侧交点，从而减弱对整体增益的依赖。这个选择也带来两个明确问题：首先，机器人需要主动完成足够宽的横切，单纯沿缆行驶未必能产生可辨识的剖面；其次，宽度对增益不敏感，并不意味着对场形不敏感。本阶段因而将搜索、检测、横切、反演和跟踪组织为连续流程，并分别讨论信号质量、模型适用性与目标身份。

在状态估计与控制方面，误差状态扩展卡尔曼滤波（error-state extended Kalman filter，ES-EKF）便于表达惯性传播、外部观测修正与协方差演化\cite{sola2017quaternion}。模型预测控制（model predictive control，MPC）则在有限预测时域内处理跟踪代价和输入约束\cite{mayne2000constrained}。将这两种方法连接起来时，主要问题是如何使用质量不稳定的观测：滤波器持续输出并不自动表示统计一致，控制器减少抖动也不必然降低横向误差。本文因此保留固定参数和不确定性调制策略的对照，按估计误差、控制变化和约束满足分别分析收益。

任务层采用行为树组织正常任务、重捕和异常处置。行为树的模块组合与优先级表达已有成熟方法基础\cite{colledanchise2018behavior}，本课题的工作重点在于把输入失效、目标不明、运动边界和执行状态接入统一的巡检流程。尤其在通信与感知异步时，“命令已经归零”和“航行器已经停住”属于不同层次的状态，必须由不同证据支持。

\section{研究目标与中期工作范围}

本课题旨在原 AUV 平台上构建多源异构感知与安全自主巡检方法。阶段目标包括获得物理含义明确的观测、向运动决策传递状态及质量，以及建立可记录、可撤回、可逐层检查的执行接口。

截至\ReportCutoff，已形成的工作包括磁搜索与主动横切反演、场形失配处理、有限巡检仿真、融合估计与控制对照，以及磁测量链、艇载计算和执行脑通信等实物子链验证。它们分别回答算法是否能够形成有效信息、策略是否具有可解释的性能变化，以及关键接口是否能够实际运行。后文将方法与代表结果相邻组织，以便把每项结论与其证据直接对应。
